\section{Results V: Real Meeting Diarization}

Meeting Mode adds a failure dimension that ordinary dictation does not have: the words may be
accurate while the speaker labels are unusable.

The campaign evaluates the complete 16-recording AMI test split, covering approximately 543.7
minutes of meeting audio and 30,714 seconds of scored reference speech. Fifteen recordings have
four reference speakers and one has three. Human RTTM labels provide the speaker reference.

\subsection{The original clustering operating point does not transfer}

% TODO(#508): generate from the AMI result artifacts.
\begin{table}[ht]
\centering
\caption{AMI test-split diarization conditions. DER is the per-recording mean shown in the
archived comparison; the unrestricted threshold-1.2 condition additionally has 27.37\%
time-weighted DER. Count error is the mean signed hypothesis-minus-reference speaker count.}
\label{tab:ami}
\begin{tabular}{lrrr}
\toprule
Configuration & Mean DER & Mean count error & Exact count \\
\midrule
Threshold 0.5, unrestricted & 75.21\% & +155.19 & 0/16 \\
Threshold 0.5, \texttt{max\_speakers=4} & 29.42\% & +0.06 & 15/16 \\
Threshold 1.2, unrestricted & 26.71\% & +2.06 & 2/16 \\
Threshold 1.2, \texttt{max\_speakers=4} & 29.42\% & +0.06 & 15/16 \\
\bottomrule
\end{tabular}
\end{table}

The former threshold of 0.5 severely over-splits the AMI meetings. The resulting
75.21\% mean DER is accompanied by a mean signed speaker-count error of +155.19. Raising the Meeting Mode clustering threshold to 1.2 reduces the
per-recording mean DER to 26.71\%. The corresponding time-weighted DER is 27.37\%, or 20.51\%
with a 250 ms collar.

The result should not be read as a claim that 26--27\% DER is state of the art. Its importance
here is diagnostic: a default calibrated on shorter or different audio was severely mismatched
to long real meetings, and real annotated evaluation changed the shipped operating point.

\subsection{A four-speaker cap nearly fixes count, not DER}

The intervention tested by the archived artifact is \texttt{max\_speakers=4}: a ceiling on
clustering, not a supplied exact speaker count. At threshold 1.2 it makes the estimated count
exact on 15 of 16 recordings, compared with 2 of 16 under unrestricted clustering. The remaining
recording has three reference speakers and is over-counted by one under the four-speaker cap.

However, the paired DER result does not establish an accuracy benefit. Seven meetings improve,
seven worsen, and two are unchanged; the exact sign test has $p=1.0$. The mean paired difference
is +2.71 DER points with a 95\% interval of -1.64 to +7.60. Thus the count constraint solves a
different problem from diarization accuracy.

\subsection{Thresholds are domain-sensitive}

The archived threshold sweeps place the preferred regions at approximately 0.8--0.9 for the
synthetic meeting fixture, approximately 0.9 for the evaluated VoxConverse condition, and
approximately 1.2 for AMI. At the AMI-oriented threshold, crowd-scene VoxConverse examples can
be under-counted badly.

This provides evidence against a single universal clustering threshold. The operating point is a
property of the embedding/clustering system together with the target domain.

\subsection{Simple centroid repair is not safe}

A follow-up experiment asks whether post-hoc cosine similarity between cluster centroids can merge
fragments belonging to the same real speaker. On the measured AMI conditions, no threshold
provides useful repair recall without introducing unacceptable merges between different people.
A very conservative threshold can achieve high precision only by fixing a small fraction of
splits; relaxing it quickly produces wrong-person merges.

We retain this as a negative result rather than tuning the heuristic until one aggregate appears
favorable.